\chapter{Dehydration}

\section*{Definition and Overview}
Dehydration occurs when the body loses more fluid than it takes in, leaving insufficient water to carry out normal functions. It can range from mild, with thirst and dry mouth, to severe, with confusion, kidney failure, and shock. Dehydration is especially dangerous in infants, young children, older adults, and people with chronic illness. Most cases respond to drinking fluids, but severe dehydration requires urgent medical treatment with intravenous fluids. This chapter covers the causes, signs, prevention, and treatment of dehydration.

\section*{Epidemiology}
Dehydration is one of the most common reasons people seek medical care and is a frequent cause of hospital admission, particularly in summer and during heatwaves. Infants and young children are highly vulnerable because of their small body size and high fluid turnover, especially with vomiting and diarrhoea. Older adults are also at risk due to reduced thirst sensation, kidney function decline, and medication effects. Gastroenteritis causes significant dehydration in children, and heat-related illness in Australia increases during hot weather.

\section*{Aetiology and Pathophysiology}
\begin{itemize}
    \item \textbf{Inadequate intake:} not drinking enough, common in older adults, people with disability, and during illness
    \item \textbf{Excess loss:} vomiting, diarrhoea, fever, sweating, and increased urination from diabetes or diuretic medicines
    \item \textbf{Heat:} hot weather and vigorous exercise increase fluid losses through sweat
    \item \textbf{Fluid shifts:} some illnesses cause fluid to move out of the bloodstream, effectively reducing circulating volume
    \item \textbf{Physiological response:} the body conserves water by reducing urine output and increasing thirst; when losses exceed intake, blood pressure falls and organs are starved of oxygen
    \item \textbf{Electrolyte disturbance:} dehydration is accompanied by imbalances of sodium, potassium, and other minerals
\end{itemize}

\section*{Clinical Features}
\begin{itemize}
    \item Thirst, dry mouth, and lips
    \item Dark, concentrated urine and reduced frequency of urination
    \item Tiredness, dizziness, and light-headedness, especially on standing
    \item Headache, dry skin, and sunken eyes
    \item In infants: reduced wet nappies, no tears when crying, a sunken soft spot (fontanelle), and irritability or listlessness
    \item Confusion, rapid heart rate, rapid breathing, and low blood pressure in severe dehydration
    \item In severe cases: fainting, reduced consciousness, and kidney failure
\end{itemize}

\section*{Differential Diagnosis}
\begin{itemize}
    \item Anaemia, which causes similar tiredness and dizziness
    \item Low blood sugar (hypoglycaemia)
    \item Heat exhaustion and heat stroke
    \item Infection with fever, which may coexist with or cause dehydration
    \item Heart or kidney failure, which can cause abnormal fluid balance
    \item Thyroid disorders and other causes of fatigue
\end{itemize}

\section*{When to Seek Medical Care}
See a doctor if someone cannot keep fluids down, is vomiting persistently, has significant diarrhoea, or shows signs of dehydration that are not improving with oral fluids. Infants, older adults, and people with chronic illness need review earlier. Medical attention is needed for confusion, persistent dizziness, or reduced urine output.

\textbf{Call 000 (or local emergency services) for:}
\begin{itemize}
    \item Severe dehydration signs: confusion, drowsiness, fainting, or very low urine output
    \item Rapid heart rate with dizziness or breathing difficulty
    \item An infant or child who is listless, unresponsive, or has not passed urine for many hours
    \item Vomiting blood or signs of shock
\end{itemize}

\section*{Diagnosis and Investigations}
\begin{itemize}
    \item \textbf{Clinical assessment:} signs such as dry mouth, reduced skin turgor, and altered consciousness indicate the severity
    \item \textbf{Urine tests:} dark colour and high concentration indicate low fluid intake; urine specific gravity may be measured
    \item \textbf{Blood tests:} electrolytes, kidney function, and haematocrit help assess severity and guide treatment
    \item \textbf{Weight measurement:} acute weight loss reflects fluid loss in children
    \item \textbf{Underlying cause:} stool tests, blood glucose, or other tests where gastroenteritis, diabetes, or infection is suspected
\end{itemize}

\section*{Management}
\begin{itemize}
    \item \textbf{Oral rehydration:} the first-line treatment; water, oral rehydration solutions, or diluted cordial in small, frequent sips
    \item \textbf{Oral rehydration solution (ORS):} the most effective for replacing fluids and electrolytes, especially in children with gastroenteritis; available from pharmacies
    \item \textbf{Intravenous fluids:} for severe dehydration, persistent vomiting, or inability to drink
    \item \textbf{Treating the cause:} antiemetics for vomiting, rehydration for diarrhoea, and management of fever or other underlying illness
    \item \textbf{Monitoring:} watching urine output, mental state, and other signs to confirm recovery
    \item \textbf{For athletes:} sports drinks and electrolyte replacement tailored to the intensity and duration of activity
\end{itemize}

\section*{Complications}
\begin{itemize}
    \item Acute kidney injury from reduced blood flow to the kidneys
    \item Electrolyte disturbances, which can cause seizures or heart rhythm problems
    \item Shock and organ failure in severe cases
    \item Heat stroke when dehydration is combined with heat exposure
    \item Constipation and urinary tract infections from chronic mild dehydration
    \item Reduced cognitive and physical performance with mild dehydration
\end{itemize}

\section*{Prognosis and Outcomes}
With appropriate treatment, dehydration resolves quickly and completely in most people. Mild to moderate dehydration responds to oral fluids within hours, and severe dehydration improves rapidly with intravenous therapy in hospital. In Australia, death from dehydration is rare when care is sought in time. The main determinant of outcome is the speed of treatment and the health of the affected person, with infants and older adults needing the most careful monitoring.

\section*{Prevention and Self-Care}
\begin{itemize}
    \item Drink regularly throughout the day, and increase intake in hot weather and during exercise
    \item Encourage infants and children to drink often, especially when unwell
    \item Offer fluids to older adults regularly, as thirst diminishes with age
    \item During illness with vomiting or diarrhoea, use oral rehydration solution and drink little and often
    \item Carry water when active, travelling, or outdoors in the heat
    \item Monitor urine colour as a simple guide to hydration
    \item Be aware that some medicines and conditions increase dehydration risk
\end{itemize}

\section*{Sources and Further Reading}
The following reputable sources provide further information for healthcare students and patients seeking reliable, current advice about dehydration.

\begin{itemize}
    \item NHS. \textit{Dehydration: symptoms, causes, and treatment.} https://www.nhs.uk/conditions/dehydration/
    \item Mayo Clinic. \textit{Dehydration: symptoms and prevention.} https://www.mayoclinic.org/diseases-conditions/dehydration/
    \item MSD Manual. \textit{Dehydration: professional overview.} https://www.msdmanuals.com/
    \item MedlinePlus. \textit{Dehydration.} https://medlineplus.gov/dehydration.html
    \item Healthdirect Australia. \textit{Dehydration and fluid replacement.} https://www.healthdirect.gov.au/
    \item World Health Organization. \textit{Oral rehydration salts: guidelines.} https://www.who.int/
\end{itemize}
